\clearpage
\subsection{Ideals}
\begin{center}
  \textcolor{gray}{------------------------------------------- 27.10.2026 --------------------------------------------} \\[0.3em]
\end{center}  

A key notion in algebra of the rings. 

\begin{definition}[title = Definition 1.4.1]
    Let $(R, +, \cdot)$ be a commutative ring. Then a subset $I$ of $R$ is called an \textit{ideal} of $R$ if the following holds: 
    \begin{itemize}
      \item[(a)] $0 \in I$ 
      \item[(b)] For all $f, g \in I$, one has $f + g \in I$ 
      \item[(c)] For all $f \in R$ and $g \in I$, one has $f \cdot g \in I$   
    \end{itemize}
\end{definition}

The definition looks a little like the notion of a vector subspace. 
% TODO: Vector subspace definition --> span of vecotrs 

You can view $I$ as a set of objects that are zero in some abstract sense. 

One can define ideals by system of generators. 

\begin{definition}[title = Definition 1.4.2]
    Let $R$ be a commutative ring and $f_1, \ldots, f_s$ finitely many elements of $R$. We define the ideal generated by $f_1, \ldots, f_s$ as follows: 
    \[\langle f_1, \ldots, f_s \rangle_{R} = \{f_1 \cdot g_1 + \ldots + f_s \cdot g_s \mid g_1, \ldots, g_s \in R\}\]  
    If the choice of $R$ is clear the subscript $R$ can be omitted.
\end{definition}

\begin{lemma}[title = Lemma 1.4.3]
    For $f_1, \ldots, f_s$ in a commutative ring $R$, the set $\langle f_1, \ldots, f_s \rangle_{R}$ is an ideal. 

    \paragraph{Proof:} 
    try poving this. Not really hard. 

    % When you wake up, brush your teeth... people that do math then prove this. 
\end{lemma}

\begin{example}[title = Example 1.4.4 (The use of ideals in the context of implicitization)]
    Let's consider an easy parametric curve given by 
      \[x = t, \quad y = t^2, \quad   z = t^3.\]
     where $t$ is a parameter. This is a space curve. Let's write those equations as 
     \[x - t = 0, \quad y - t^2 = 0, \quad z - t^3 = 0, \] 
    where $x-t, y-t^2, z-t^3 \in k[x_1, \dots, x_n]$. The ideal $I = \langle x-t, y-t^2, z-t^3 \rangle$ contains polynomials $f(x,y,z,t)$ that define equations $f(x,y,z,t) = 0$
    deriviable out of our thee equations (derivable in the sense of 1.4.2).   

    It turns out that some polynomials in $I$ don't depent on $t$. 
    \begin{align*}
        \left. 
        \begin{aligned}
            (x+t) \cdot (x-t) &+ (-1)\cdot (y-t^2) = x^2 -y \\
            (x^2 + xt + t^2) \cdot (x-t) &+ (-1) \cdot (z - t^3)  = x^3 - z 
        \end{aligned} \right\} \leftarrow \text{ gives two implicit equations for our curve.}
    \end{align*} 
    In more complicated situations, you won't see the implicit equations to easily.   

    \begin{minipage}{0.59\textwidth}
       \begin{center}
        \begin{tikzpicture}[x=0.01cm, y=-0.01cm]
          \draw[blue!75, thick, use Hobby shortcut]
            (5,168) .. (65,60) .. (200,15) .. (300,45) .. (335,120)
            .. (290,220) .. (180,293) .. (110,230) .. (140,150)
            .. (250,110) .. (350,155) .. (400,225) .. (520,273)
            .. (650,245) .. (748,160) .. (700,65) .. (620,55)
            .. (545,115) .. (505,250) .. (450,375) .. (295,447)
            .. (180,400) .. (145,270) .. (190,110) .. (290,25)
            .. (420,32) .. (550,115) .. (620,205) .. (667,375);
        \end{tikzpicture} 
        \end{center} 
    \end{minipage}
    \hspace*{0.2cm}
    \begin{minipage}{0.39\textwidth}
        \begin{align*}
            x &= t - t^2 + 3t^3 - t^{10} \\
            y &= t + t^5  - t^7
        \end{align*}
    \end{minipage}  
\end{example} 

\clearpage
\begin{example}[title = Example 1.4.4 Continued]
  
There is a maschinery that helps us to find implicit equations in such cases. It is called \textit{Gröbner basis}. We can use \textit{SageMath} to solve such problems.\\ 

\begin{sagecode}{SageMath Code}
R = PolynomialRing(QQ, 'x,y,t') 
x,y,t = R.gens() 
I = R.ideal(x -(t - t^2 + 3*t^3 - t^10), y - (t + t^5 - t^7)) 
I.elimination_ideal(t) <-- done with Gröbner bases
\end{sagecode} 

Let $u(t),v(t) \in k{t}$ and a parametric curve given by    
\[x = u(t), \quad y = v(t)\] 
Introduce the ideal $\langle x - u(t), g - u(t) \rangle \subseteq k[x,y,t]$, one has the so-called \textit{eliminaton ideal} $I \cap k[x,y]$. The task is to generate this ideal. 
\end{example}


\begin{example}[title = Example 1.4.4a] 
    \begin{minipage}[t]{0.49\textwidth}
        \begin{center}
            \begin{tikzpicture}[x=1cm, y=0.5cm]

                \draw[->] (-0.5,0) -- (5.5,0) node[right] {$x$};
                \draw[->] (0,-3) -- (0,6) node[above] {$y$};

                \draw[thick, kGreen, domain=-0.2:5.2, samples=60, smooth, variable=\x]
                    plot ({\x}, {\x*\x - 5*\x + 4});
            \end{tikzpicture}
        \end{center}        
        Discriminant $\neq$ 1
    \end{minipage}
    \hspace*{0.2cm}
    \begin{minipage}[t]{0.49\textwidth}
         \begin{center}
            \begin{tikzpicture}[x=1cm, y=0.5cm]

                \draw[->] (-0.5,0) -- (5.5,0) node[right] {$x$};
                \draw[->] (0,-3) -- (0,6) node[above] {$y$};

                \draw[thick, kGreen, domain=-0.2:4.3, samples=60, smooth, variable=\x]
                  plot ({\x}, {\x*\x - 4*\x + 4}); 

                \fill[kGreen] (2,0) circle (2pt);
            \end{tikzpicture}
        \end{center}  
        Discriminant = 0 \\
        Derivative = 0 
         \end{minipage}
         \[I = \langle x^2 + px + q, 2x+p \rangle \subseteq k[x, p, q]\] 
        \[I \cap k[p, q] \leftarrow \text{ the generator will essentially be the discriminant}\]
   
\end{example}

\begin{definition}[title = Definition 1.4.5]
    An ideal $I$ in a commutative ring $R$ is called \textit{finitely generated} if $I = \langle f_1, \ldots, f_s \rangle$ for some finitely many $f_1, \ldots, f_s \in R$. 
    In this case $\{f_1, \ldots, f_s\}$ is called a \textit{set of generators} of $I$.
\end{definition}

\begin{tcolorbox}[sharp corners, colframe=black!75!black]
    \textbf{Key Question:} Are ideals in the $k[x_1, \dots, x_n]$ \underline{finitely} generated?  
\end{tcolorbox}

Now, we have an interplay: ideals vs. varieties. 
\[V(f_1, \ldots, f_s) \quad \quad \quad \langle f_1, \ldots, f_s \rangle\]
One can view the variety $V(f_1, \ldots, f_s)$ not as a function of $(f_1, \ldots, f_s)$ but as a function og the ideal $\langle f_1, \ldots, f_s, \rangle$. 

\clearpage
\begin{proposition}[title = Proposition 1.4.6]
    % Let $f_1, \ldots, f_s, g_1_, \ldots, g_s \in k[x_1, \dots, x_n]$ be such that 
    \[\langle f_1, \ldots, f_s \rangle = \langle g_1, \ldots, g_s \rangle\] 
    Then \[V(f_1, \ldots, f_s) = V(g_1, \ldots, g_s).\] 

    \paragraph{Proof:} You'll derive it by spelling out the underlying notation. Try! 
\end{proposition} 

\begin{definition}[title = Definition 1.4.7]
    For a set $X \subseteq k^n$ we introduce the \textit{ideal of $X$} by: 
    \[I(X) = \{f \in k[x_1, \dots, x_n] \mid f(a) = 0 \text{ for all } a \in X\}\]
\end{definition}

\begin{lemma}[title = Lemma 1.4.8]
    For $X \subseteq k^n$, the set $I(X)$ is in fact an ideal. 

    \paragraph{Proof:} Just spell out $I(X)$ and use the definition of an ideal. Try! 
\end{lemma}

\begin{example}[title = Example 1.4.9]
    \begin{enumerate}
      \item $V = \{(0,0)\}$ variety in $k^2$. \[I(V) = \{f \in k[x,y] \mid f(0,0) = 0\} = \langle x, y \rangle\] 
      Polynomials in the two variabeles: \[\sum_{i, j} c_{ij} x^i y^j \quad \Rightarrow f(0,0) = c_{0,0}\]
    
    \end{enumerate}
\end{example}